\scanpage{8}
\mspage{4}
une application
\[
(22)\qquad\widetilde{\mathfrak X}'(i)\xrightarrow{\Phi_i}\operatorname{Hom}(\mathbf T(\overline{\mathbf Q}),\pi_1(\mathfrak X'))
\]
\margin{ici il faut supposer $\dim\mathfrak X=1$ -- autrement il
faut \unsure{voir|revoir} la définition de $\widetilde{\mathfrak X}'(I)$.....\\
autre façon de dire -- il faut que $\widetilde{\mathfrak X}'$ \unsure{soit considéré}
\unsure{pour donner des|pour définir des} \unsure{traits} \unsure{henséliens ou locaux au}
dessus d'un pt de $\unsure{\overline{\mathfrak X}'(I)|\mathfrak X'(I)}$\unsure{.....|....}}
compatible avec les opérations de $\pi_1(\mathfrak X')$
de part et d'autre, et qui précise
la donnée de $\unsure{\varphi|\varphi_i}$ (17). Mieux encore,
le système $(\Phi_i)_{i\in I}=\Phi$.
\[
(23)\qquad\widetilde{\mathfrak X}'(I)\xrightarrow{\Phi}\operatorname{Hom}(\mathbf T(\overline{\mathbf Q}),\pi_1(\mathfrak X'))
\]
est compatible avec les $\unsure{\mathcal E|\Sigma}_{\mathfrak X'}\deleted{\unsure{\mathcal E}}$ de part et
d'autre, en se rappelant, pour l'opération
de $\unsure{\mathcal E|\Sigma}_{\mathfrak X'}$ sur le mb de droite, de tenir
compte aussi de l'opération de $\unsure{\mathcal E|\Sigma}_{\mathfrak X'}$
sur $\mathbf T(\overline{\mathbf Q})$ (via $\Gamma$).

\textcircled{5} Soit enfin, par dessus le marché, une
partie $J\subset\mathfrak X'(\overline{\mathbf Q})$. Soit $\widetilde{\mathfrak X}'(J)$ son
image inverse dans $\widetilde{\mathfrak X}'$. Soit $\mathcal G_{\mathfrak X'}(I,J)$
ou simplement $\mathcal G'$ le sous-groupe de $\mathcal G_{\mathfrak X'}(I)$
qui invarie $J$. \deleted{D. Voici les structures} \inserted{Donc $\mathcal G'$}
opère sur $J$, sur $\widetilde{\mathfrak X}'(J)$ et sur
\[
(24)\qquad\widetilde{\mathfrak X}'(J)\longrightarrow J
\]
\mspage{7}
